% ============================================================
% paperflow article template — single-column SCI journal style
% All double-at placeholders are filled by the paperflow pipeline
% via plain Python string replacement (contract: see README.md).
% Do NOT rename them.
% Compiles with: tectonic template_filled.tex
% Only TeX Live core packages are used (no shell-escape, no minted).
% ============================================================
\documentclass[11pt,a4paper]{article}

% ---------- page layout ----------
\usepackage[margin=2.5cm]{geometry}

% ---------- math (load amsmath BEFORE newtxmath) ----------
\usepackage{amsmath}

% ---------- Times-like text and math fonts ----------
\usepackage{newtxtext}
\usepackage{newtxmath}

% ---------- graphics and tables ----------
\usepackage{graphicx}
\usepackage{booktabs}

% ---------- captions ----------
\usepackage[labelfont=bf,font=small]{caption}

% ---------- citations: author-year, (Author, 2024; Other, 2025) ----------
\usepackage[round,semicolon,authoryear]{natbib}

% ---------- line spacing ----------
\usepackage{setspace}
\setstretch{1.15}

% ---------- hyperref MUST be loaded last ----------
\usepackage[hidelinks]{hyperref}

% ---------- title block ----------
\title{Iron Incorporation Turns Nickel Hydroxide into a High-Activity Oxygen Evolution Electrocatalyst: An Illustrative NiFe-LDH Case Study}
\author{paperflow demo\\ \small academic-skills / paperflow pipeline demonstration}
\date{\today}

\begin{document}

\maketitle

\begin{abstract}
The oxygen evolution reaction (OER) limits the efficiency of alkaline water electrolysis, and inexpensive catalysts that rival noble-metal oxides are therefore in continuous demand \citep{mccrory2013}. In this demonstration study we examine, on an illustrative example dataset, how iron incorporation transforms nickel hydroxide into a high-activity NiFe layered double hydroxide (NiFe-LDH) electrocatalyst. Powder X-ray diffraction confirms a phase-pure hydrotalcite-type structure, while Ni 2p X-ray photoelectron spectroscopy resolves coexisting Ni2+ and Ni3+ states consistent with Fe-modulated electronic structure \citep{trotochaud2014}. In 1 M KOH, the illustrative polarization data give an overpotential of 273 mV at 10 mA cm-2 for NiFe-LDH, surpassing both the RuO2 benchmark (324 mV) and Fe-free Ni(OH)2 (466 mV), and the accompanying Tafel slope of 38 mV per decade indicates fast electron-transfer kinetics \citep{dionigi2016}. These trends reproduce the well-established activity ordering of the NiFe (oxy)hydroxide family and illustrate an end-to-end, citation-verified manuscript-generation workflow in which every quantitative claim is traceable either to the declared example dataset or to literature that survived authenticity checking. The workflow, rather than the chemistry, is the deliverable of this document, and all numbers herein are illustrative rather than experimental.
\end{abstract}

\medskip
\noindent\textbf{Keywords:} oxygen evolution reaction, NiFe-LDH, electrocatalysis, workflow automation

\begin{center}\fbox{\parbox{0.9\linewidth}{\small\itshape DEMO ONLY - generated by the paperflow pipeline with an illustrative synthetic dataset; not real experimental results and not for submission.}}\end{center}

\section{Introduction}

Electrochemical water splitting stores renewable electricity as hydrogen, but its overall efficiency is throttled by the sluggish four-electron oxygen evolution reaction at the anode. Benchmarking studies have standardized the comparison of OER catalysts through the overpotential required to reach 10 mA cm-2, a metric tied to the photocurrent of an ideal solar-driven device \citep{mccrory2013}. Within design-principle frameworks that correlate activity with electronic-structure descriptors, transition-metal oxides and (oxy)hydroxides have emerged as the most promising earth-abundant candidates in alkaline media \citep{suntivich2011,grimaud2013}. Among them, nickel-iron systems occupy a special position: seminal work demonstrated that even incidental iron impurities dramatically enhance the apparent activity of nickel (oxy)hydroxide, identifying Fe incorporation as the true origin of its celebrated performance \citep{trotochaud2014}. Subsequent studies of NiFe layered double hydroxides consolidated this family as the state of the art for non-acidic OER, with overpotentials frequently below 300 mV \citep{dionigi2016}. The present manuscript revisits this well-charted chemistry for a different purpose. It serves as an end-to-end demonstration of paperflow, a pipeline that drafts, cites, verifies, and assembles a manuscript automatically, using an illustrative NiFe-LDH dataset. Our aim is to show that each stage, from outline to citation-authenticity gating to document assembly, can be orchestrated reproducibly while keeping every factual claim anchored to verified sources.

\section{Results and Discussion}

Figure 1 summarizes the illustrative structural and electrochemical dataset. The diffraction pattern of NiFe-LDH in Figure 1a reproduces the characteristic (003), (006), (012), and (015) reflections of a hydrotalcite-type layered structure and matches the LDH reference card, whereas the Fe-free control exhibits the distinct beta-Ni(OH)2 pattern; no impurity phases are apparent in either trace. The Ni 2p spectrum in Figure 1b is deconvoluted into Ni2+ and Ni3+ components together with their shake-up satellites, and the appreciable Ni3+ fraction is consistent with the Fe-induced modulation of nickel electronic structure that has been linked mechanistically to enhanced OER activity in mixed Ni-Fe (oxy)hydroxides \citep{trotochaud2014}. Turning to catalysis, the polarization curves in Figure 1c, acquired on the illustrative dataset in 1 M KOH, place the three electrodes in the order NiFe-LDH > RuO2 > Ni(OH)2, with overpotentials at 10 mA cm-2 of 273, 324, and 466 mV, respectively, following the benchmarking convention of the field \citep{mccrory2013}. The Tafel analysis in Figure 1d yields slopes of 38, 42, and 65 mV per decade for the same ordering, indicating that Fe incorporation accelerates the rate-determining electron-transfer step rather than merely increasing the electrochemically active area. Both the absolute overpotential range and the roughly one-hundred-millivolt advantage over the nickel-only control agree with the consensus picture assembled for the NiFe (oxy)hydroxide family \citep{dionigi2016}, lending face validity to the synthetic dataset even though it is illustrative rather than experimental.

\begin{figure}[htbp]
\centering
\includegraphics[width=0.92\linewidth]{figures/fig_oer_generated.png}
\caption{Structural and electrochemical characterization of NiFe-LDH (illustrative synthetic data): (a) XRD patterns with the LDH reference; (b) Ni 2p XPS with peak deconvolution; (c) OER polarization curves with the eta10 annotation; (d) Tafel analysis.}
\label{fig:oer}
\end{figure}

\section{Conclusion}

This demonstration assembled a complete, citation-verified manuscript around an illustrative NiFe-LDH oxygen-evolution dataset. The narrative reproduces the accepted understanding of the system: a phase-pure layered double hydroxide, mixed Ni2+/Ni3+ surface states, and an alkaline OER activity ranking of NiFe-LDH ahead of RuO2 and far ahead of Fe-free Ni(OH)2, in line with the benchmarking literature and mechanistic studies of iron incorporation \citep{mccrory2013,trotochaud2014,dionigi2016}. More importantly for the purpose of this document, every reference cited here passed an automated authenticity gate that queries Crossref and Semantic Scholar and drops any candidate whose DOI or title cannot be confirmed; one deliberately fabricated reference planted in the candidate list was rejected by this gate and therefore does not appear in the bibliography. The pipeline thus enforces, by construction, the separation between verifiable literature claims and the clearly labeled illustrative dataset. Future work on the pipeline itself includes journal-specific formatting profiles, richer figure-placement logic, and a reviewer-simulation stage, while any scientific use of the chemistry discussed here would of course require real experimental data in place of the example dataset used throughout this demonstration.


\bibliographystyle{plainnat}
\bibliography{references}

\end{document}
